
\documentclass[a4paper]{article}

\usepackage[a4paper, top=8mm, bottom=8mm, left=8mm, right=8mm]{geometry}

\usepackage{polyglossia}
\setdefaultlanguage[babelshorthands=true]{russian}
\setotherlanguage{english}

\usepackage{fontspec}
\setmainfont{FreeSerif}
\newfontfamily{\russianfonttt}[Scale=0.7]{DejaVuSansMono}

\usepackage[tiny, compact]{titlesec}

\usepackage{titling}
\setlength{\droptitle}{-1cm}
\pretitle{\begin{center}\begin{bfseries}\Large}
\posttitle{\par\end{bfseries}\end{center}}
\preauthor{\begin{center}\normalsize}
\postauthor{\par\end{center}\vspace{-1.8cm}}

\usepackage{hyperref}
\usepackage{bookmark}
\usepackage{csquotes}
\usepackage{amsmath}

\title{Реализация свёртки изображений на Python}

\author{Федотов Дмитрий Алексеевич}

\date{}

\begin{document}

\maketitle

\begin{flushright}
    Группа: \emph{\enquote{25.Б41-мм}}

    Кафедра: \emph{системного программирования}

    Научный руководитель: \emph{Пономарев Николай Алексеевич}

    Номер семестра практики: \emph{2}
\end{flushright}

\section{Постановка задачи}

Свёртка является одной из базовых операций цифровой обработки изображений: с её помощью выполняются размытие, повышение резкости и выделение границ. В готовых библиотеках эта операция обычно скрыта за высокоуровневыми функциями, поэтому для понимания алгоритма полезно реализовать её самостоятельно и сравнить результат с промышленными решениями.

Работа выполнялась в рамках учебного проекта по цифровой обработке изображений. \textbf{Целью работы является} разработка собственной реализации операции свёртки изображений на языке Python, а также проверка корректности реализации и сравнение её производительности с готовыми библиотечными решениями.

Результаты работы полезны как учебная реализация базового алгоритма компьютерного зрения: проект позволяет явно увидеть, как ядро свёртки применяется к пикселям изображения, как обрабатываются граничные пиксели и почему специализированные библиотеки оказываются быстрее собственной реализации.

Для достижения цели были поставлены следующие задачи:
\begin{itemize}
    \item реализовать свёртку для grayscale- и RGB-изображений;
    \item поддержать ядра размера \(3 \times 3\) и \(5 \times 5\);
    \item реализовать режимы обработки краёв \texttt{reflect}, \texttt{replicate} и \texttt{wrap};
    \item добавить набор ядер для размытия, повышения резкости и выделения границ;
    \item реализовать CLI-интерфейс для выбора изображения, фильтра, режима обработки края и цветового режима;
    \item проверить корректность реализации тестами и провести бенчмарки относительно OpenCV и Pillow.
\end{itemize}

\section{Описание предлагаемого решения}

Свёртка изображения описывается формулой
\[
    C[i,j] = \sum_{k=-p}^{p}\sum_{l=-p}^{p} A[i+k,j+l]B[p+k,p+l],
\]
где \(A\) --- исходное изображение, \(B\) --- ядро размера \((2p+1) \times (2p+1)\), а \(C\) --- результат фильтрации.

Реализация выполнена на Python с использованием NumPy и Pillow. NumPy применяется для представления изображений в виде массивов и для векторизованных вычислений, Pillow --- для чтения изображений и преобразования результата обратно в изображение. В проекте также используется pytest для функциональных тестов и pytest-benchmark для измерения производительности.

Основная функция \texttt{convolution} принимает изображение, ядро свёртки и функцию обработки индексов за пределами изображения. Для grayscale-изображений временно добавляется ось каналов, чтобы общий алгоритм одинаково работал и с двумерными, и с RGB-изображениями. После вычисления результат ограничивается диапазоном \([0, 255]\) и приводится к типу \texttt{uint8}.

В проекте реализованы следующие группы ядер:
\begin{itemize}
    \item размытие: \texttt{box\_blur}, \texttt{gaussian\_blur};
    \item выделение границ: \texttt{sobel\_x}, \texttt{prewitt\_y}, \texttt{laplacian}, \texttt{sobel\_x\_5}, \texttt{sobel\_y\_5}, \texttt{laplacian\_5};
    \item повышение резкости: \texttt{sharpen}, \texttt{sharpen\_5}.
\end{itemize}

Для обработки краёв изображения реализованы три режима. Режим \texttt{reflect} отражает координату относительно края, \texttt{replicate} использует ближайший допустимый пиксель, а \texttt{wrap} циклически переносит индекс на противоположную сторону изображения.

\section{Апробация / Эксперименты}

Корректность реализации проверялась тестами с эталонными изображениями. Для каждого реализованного ядра результат собственной функции сравнивается с заранее подготовленным эталоном. Поскольку Pillow не предоставляет явного выбора режима обработки краёв для используемой операции фильтрации, сравнение выполняется по внутренней области изображения без краевой рамки шириной в половину размера ядра. Это позволяет проверить основное вычисление свёртки независимо от различий в приграничных пикселях.

Производительность измерялась с помощью pytest-benchmark. В эксперименте сравнивались три реализации: собственная реализация, OpenCV \texttt{cv2.filter2D} и фильтрация через Pillow. Измерения проводились на случайно сгенерированных изображениях размеров \(128 \times 128\), \(256 \times 256\), \(512 \times 512\), \(1024 \times 1024\), \(1920 \times 1080\), \(2560 \times 1440\), \(3840 \times 2160\) для grayscale- и RGB-режимов.

В выборочных результатах бенчмарка указывались среднее время и стандартное отклонение в формате \(\mathrm{mean} \pm \mathrm{stddev}\). Например, для изображения \(512 \times 512\) в режиме \texttt{reflect} собственная реализация выполняла \texttt{box\_blur} за \(41.33 \pm 0.58\) мс на grayscale-изображении и за \(96.64 \pm 1.25\) мс на RGB-изображении. Для тех же случаев OpenCV показывал \(0.210 \pm 0.005\) мс и \(0.598 \pm 0.011\) мс соответственно, а Pillow --- \(0.751 \pm 0.034\) мс и \(2.75 \pm 0.06\) мс.

По результатам экспериментов OpenCV оказался быстрее собственной реализации в среднем примерно в 188 раз, а Pillow --- примерно в 23 раза. При этом собственная реализация масштабируется предсказуемо: при увеличении количества пикселей время работы растёт близко к линейной зависимости от размера изображения.

Основные угрозы корректности эксперимента связаны с тем, что измерения проводились на одной аппаратной конфигурации, а входные изображения для бенчмарков генерировались случайно и не включали затраты на чтение с диска. Кроме того, разные библиотеки по-разному обрабатывают граничные пиксели и округление значений, поэтому корректность проверялась отдельно от приграничной области.

\section{Заключение}

В ходе выполнения работы получены следующие результаты:

\begin{itemize}
    \item реализована операция свёртки изображений на Python с поддержкой grayscale- и RGB-изображений;
    \item реализованы ядра для размытия, повышения резкости и выделения границ;
    \item реализованы режимы обработки краёв \texttt{reflect}, \texttt{replicate} и \texttt{wrap};
    \item добавлен CLI-интерфейс для выбора входного изображения, фильтра и параметров обработки;
    \item подготовлены тесты корректности по эталонным изображениям;
    \item проведено сравнение производительности собственной реализации с OpenCV и Pillow.
\end{itemize}

Материалы, иллюстрирующие выполнение работы, размещены в репозитории проекта: \url{https://github.com/Fedotov-Dmitriy/image-convultion-2sem}.

Репозиторий/ссылка на пуллреквест: \url{https://github.com/Fedotov-Dmitriy/image-convultion-2sem}.

Техническая документация представлена в файле README репозитория и в файле с анализом бенчмарков \texttt{benchmark\_analysis.md}.

Решение находится в завершённом состоянии для учебной практики: реализованы основные функции, подготовлены тесты и воспроизводимые бенчмарки.

\end{document}
